\subsection{Malle's character permutations and the technical lemma}
\label{Section: Malle's character permutations and the technical lemma}

The fake degree $P_{\chi}(q)\coloneqq \sum q^{e_i(\chi)}$ of an irreducible character $\chi\in\widehat{W}$ is a polynomial that records the exponents $e_i(\chi)$ of the character (see \cite[\S4.4]{lehrer-taylor-book}). Beynon and Lusztig \cite[Prop.~A]{beynon-lustzig-coxeter-numbers-genesis} had observed a remarkable reciprocity property for these polynomials for Weyl groups. They satisfy $$P_{\chi}(q)=q^{c_{\chi}}P_{\iota(\chi)}(q^{-1}),$$where $c_{\chi}$ is the Coxeter number\footnote{However, Beynon and Lusztig, and later Malle, did not assign an epithet for these numbers; the mathematical godfathers were Gordon and Griffeth \cite{gordon-griffeth-catalan-numbers} who named them after Coxeter.} as given in Defn.~\ref{Defn: Coxeter numbers} and $\iota$ is a permutation of the irreducible characters that for Weyl groups is the identity apart from two characters of $E_7$ and four of $E_8$. 

Malle later on \cite[Thm.~6.5]{malle-fake-degrees} extended this reciprocity result to all complex reflection groups, defining a permutation $\Psi$ of the characters that is induced by a Galois action on the irreducible characters of the Hecke algebra (the two permutations satisfy $\iota(\chi)=\Psi(\chi^*)$). This permutation of Malle  is exactly the missing ingredient for the proof of Lemma~\ref{Lemma: only multiples of |g| contribute}; the characters $\chi$ for which $c_{\chi}$ is not a multiple of $|g|$ are grouped together by $\Psi$ and their contributions cancel.

\subsubsection*{A Galois action on the characters}

Recall (see \eqref{EQ: admissible specializations} and \eqref{EQ: admissible specializations and Tits factoring}) the specializations of the Hecke algebra $\mathcal{H}_{\bm x}(W)$ and $\mathcal{H}_x(W)$ that have coefficient fields $K(\bm x)$ and $K(x)$, and splitting fields $K(\bm y)$ and $K(y)$ respectively. Recall also that, after Prop.~\ref{Prop. Malle's charact. of splitting field L of H} the parameters satisfy $y_{\mathcal{C}}^{N_W}=x_{\mathcal{C}}$ and $y^{N_W}=x$.

\begin{defi}\label{Defn: Malle's character permutations}
We consider the permutations $\Psi_{\mathcal{C}}$ and $\Psi$ acting on the sets $\op{Irr}(\mathcal{H}_{\bm x}(W))$ and $\op{Irr}(\mathcal{H}_x(W))$ that are respectively induced by the Galois automorphisms $\Sigma_{\mathcal{C}}$ (for $\mathcal{C}\in\mathcal{A}/W$) and $\Sigma$:\begin{align*}
\Sigma_{\mathcal{C}}&\in\op{Gal}\big(K(\bm y)/K(\bm x)\big) &\Sigma &\in\op{Gal}\big(K(y)/K(x)\big)\\
y_{\mathcal{C}}&\mapsto e^{2\pi i/N_W}\cdot y_{\mathcal{C}} & y &\mapsto e^{2\pi i/N_W}\cdot y
\end{align*} 
In particular, they are defined via $\Psi_{\mathcal{C}}(\chi_{\bm y})(T_{\bm g})\coloneqq \Sigma_{\mathcal{C}}\big( \chi_{\bm y}(T_{\bm g})\big)$ and similarly for $\Psi$. By Tits' deformation theorem, they induce permutations on the set $\widehat{W}$ of irreducible characters of $W$, which we also denote by $\Psi_{\mathcal{C}}$ and $\Psi$.
\end{defi}

The permutations $\Psi_{\mathcal{C}}$ and $\Psi$ satisfy a set of properties with respect to the Coxeter numbers and other statistics of the characters $\chi\in\widehat{W}$:

\begin{prop}\label{Prop: properties of Psi for Cox numbers}
For any character $\chi\in\widehat{W}$ and orbits $\mathcal{C},\mathcal{C}'\in\mathcal{A}/W$, the following are true:
\begin{tasks}[style=enumerate, after-item-skip=4mm](3)
\task $\displaystyle\Psi_{\mathcal{C}}(\chi)(1)=\chi(1)$\label{task1}
\task $\displaystyle m^{\Psi_{\mathcal{C}}(\chi)}_{\mathcal{C}',j}=m^{\chi}_{\mathcal{C}',j}$\label{task2}
\task $\displaystyle c_{\Psi_{\mathcal{C}}(\chi)_{,\mathcal{C}'}}=c_{\chi_{,\mathcal{C}'}}$\label{task3}
\end{tasks}
\end{prop}
\begin{proof}
Since $\Psi_{\mathcal{C}}$ is induced by a Galois automorphism, it has to respect the degree of the character $\chi_{\bm y}$, hence also of $\chi$; this proves part~\ref{task1} The spectrum of any braid reflection $\bm s_{\mathcal{C}',\gamma}$ is generically rational (see Defn.~\ref{Defn: rational character}) by the defining relations \eqref{EQ: deformed order relations}. This means that the eigenvalues of any $\bm s_{\mathcal{C}',\gamma}$ in the representation that affords $\chi_{\bm y}$ live in the coefficient field $K(\bm x)$ and are therefore fixed by $\Psi_{\mathcal{C}}$. This proves part~\ref{task2} after recalling the definition of $m^{\chi}_{\mathcal{C},j}$ from the start of \S~\ref{Section: Character values on roots of the full twist} and also part~\ref{task3} after Prop.~\ref{Prop: local cox num formula}. The same results are of course true for $\Psi$.
\end{proof}

The following is the key technical lemma that we have been building towards through all of \S~\ref{Section: Hecke algebras and the technical lemma}. The character calculations of Prop.~\ref{Prop: character values on roots of full twist} were included just so that the argument presented here is self-contained.

\begin{prop}[The key technical lemma]\label{Prop: The key technical lemma}
Let $g$ be a $\zeta$-regular element of $W$, $\chi\in\widehat{W}$ an irreducible character, and $\mathcal{C}\in\mathcal{A}/W$ an orbit of hyperplanes. Then, we have $$\Psi_{\mathcal{C}}(\chi)(g)= \zeta^{-c_{\chi_{,\mathcal{C}}}}\cdot\chi(g) \quad \quad\text{and}\quad\quad\Psi(\chi)(g)= \zeta^{-c_{\chi}}\cdot \chi(g).$$
\end{prop}
\begin{proof}
Assume that $\zeta=e^{2\pi il/d}$ with $(l,d)=1$. Then, by Prop.~\ref{Prop: regular elements lift to roots of the full twist}, we can lift  $g$ to some element $\bm g\in B(W)$ that is commensurable with the full twist (i.e. it satisfies $\bm g^d=\bm \pi^l$). Now, replacing $x_\mathcal{C}$ with $y_{\mathcal{C}}^{N_W}$ we can rewrite the character evaluations from \eqref{EQ: character values via local Coxeter} as $$\chi_{\bm y}(T_{\bm g})=\chi(g)\cdot\prod_{\mathcal{C}'\in\mathcal{A}/W}y_{\mathcal{C}'}^{N_W(e_{\mathcal{C}'}\omega_{\mathcal{C}'}-c_{\chi,_{\mathcal{C}'}})l/d}, $$
which, after applying the Galois automorphism $\Sigma_{\mathcal{C}}$, becomes
$$ \Psi_{\mathcal{C}}(\chi_{\bm y})(T_{\bm g})= \chi(g)\cdot e^{2\pi i(e_{\mathcal{C}}\omega_{\mathcal{C}}-c_{\chi_{,\mathcal{C}}})l/d }\cdot \prod_{\mathcal{C}'\in\mathcal{A}/W}y_{\mathcal{C}'}^{N_W(e_{\mathcal{C}'}\omega_{\mathcal{C}'}-c_{\chi,_{\mathcal{C}'}})l/d}. $$ 
Now, this is really 
$$ \Psi_{\mathcal{C}}(\chi_{\bm y})(T_{\bm g})=e^{2\pi i(e_{\mathcal{C}}\omega_{\mathcal{C}}-c_{\chi_{,\mathcal{C}}})l/d }\cdot\chi_{\bm y}(T_{\bm g}),$$ which completes the proof after applying Tits' deformation theorem and recalling that $e_{\mathcal{C}}\omega_{\mathcal{C}}$ is a multiple of $d$ by Corol.~\ref{Corol: regular numbers divide e_co_c}. The same argument of course works for~$\Psi$.
\end{proof}
